\section{Introduction}

Evaluation for intelligent vehicles has long relied on simulation and scenario-based testing because exhaustive real-world validation is infeasible, expensive, and slow to reproduce \cite{riedmaier2020scenario}. That challenge becomes sharper when the decision-making component is a large language model. Recent work has shown that LLMs can serve as planners, controllers, or interpretable decision interfaces in autonomous-driving settings \cite{wen2023dilu,sha2023languagempc,fu2023drivelikehuman,xu2023drivegpt4}. However, once these agents are evaluated under local deployment constraints, task quality is no longer the whole story. Fixed inference budgets, timeout escalation, and fallback behavior can change whether a run remains comparable at all.

This issue matters directly for automated-vehicle evaluation practice. In simulation studies, benchmark rankings are often used to choose controller variants, adaptation strategies, or the next round of validation effort. If invalid runs are silently filtered or interpreted as ordinary score outliers, the evaluation can recommend models that do not remain operational under the runtime regime in which they are actually deployed. Prior work on consumer-grade LLM serving already shows that practical usability depends heavily on hardware-aware runtime behavior \cite{song2023powerinfer}. In local intelligent-vehicle evaluation, that same runtime behavior determines which empirical conclusions survive scrutiny.

Fine-tuning makes this measurement problem more important rather than less. Parameter-efficient methods such as QLoRA make local adaptation feasible on limited hardware \cite{dettmers2023qlora}. But an adapted model is only meaningfully better if its gain survives the runtime regime used to measure it. Broader model evaluation work has already argued that trustworthy assessment requires more than a single aggregate score \cite{liang2022helm}. In our setting, that principle takes a concrete form: if fine-tuning improves action quality but pushes a local driving agent into timeout collapse, the correct conclusion is not a simple fine-tuning win, but a tradeoff between policy quality and benchmark validity.

DiLu-Ollama addresses this problem by combining a DiLu-style driving decision loop with local Ollama inference, \metric{highway-env} simulation, a task-conditioned highway benchmark, and a post-hoc study pipeline that preserves invalidity and failure context rather than hiding it \cite{wen2023dilu,leurent2018highwayenv}. We use that pipeline to study whether fine-tuning improves small-to-mid local driving models under a shared three-tier highway benchmark. The evidence is deliberately narrow: highway only, no cross-scenario empirical claims, and pairwise fine-tuning conclusions restricted to exact ranking-eligible base-versus-tuned pairs. Under those rules, the usable evidence is concentrated in lightweight, where the Phi family yields one exact eligible improvement. Midclass remains screening-only, and highclass becomes a meaningful result of its own: the point where the benchmark regime is fragile enough that competitive ranking claims are no longer justified.

This paper makes three contributions. First, it presents a local-first evaluation and analysis scaffold for intelligent-vehicle LLM agents in which benchmark validity is part of the empirical outcome. Second, it provides a three-tier study showing that benchmark viability degrades with model tier under a fixed local runtime policy. Third, it provides initial but concrete evidence that fine-tuning can improve local driving agents in lightweight settings while making explicit where broader family-level claims are still blocked.
